\section{Introduction}
\label{intro}

Online recruitment fraud is the posting of fake job advertisements, usually to collect money or
personal data from applicants~\cite{vidros2017}. A job board cannot have every ad read by a
moderator, so it needs a model that ranks incoming ads by how likely they are to be fake. We study
this task on EMSCAD~\cite{vidros2017}, 17,880 job ads labelled real or fraudulent by hand. Two
properties of the data make it harder than it looks: only 4.84\% of the ads are fraudulent, so a
model that calls every ad real already reaches 95.2\% accuracy, and many ads share their text, so a
random split puts copies of the same ad in both training and test rows. Published results range
from 91\% to 99.9\% accuracy, but under different splits, metrics and resampling schemes, and
usually on a single split~\cite{taneja2025,amaar2022}.

We therefore compare models under one protocol that does not leak and that measures the spread
over several splits. The study centres on three models that see the problem differently: logistic
regression on TF-IDF features, the standard text baseline; a one-class SVM, which learns only what
real ads look like and needs no fraud labels; and a fine-tuned DistilBERT transformer. Since
DistilBERT and logistic regression make different mistakes, we also combine them by stacking.
The stack is the best model (test PR-AUC 0.867), while logistic regression (0.821) and DistilBERT
(0.816) are level and the one-class SVM reaches 0.176, above the 0.049 floor.

Section~\ref{related} reviews earlier work, Section~\ref{data} explores and prepares the data,
Section~\ref{method} describes the protocol and the models, and Sections~\ref{results}
to~\ref{conclusion} present, explain and discuss the results. The code is organised as four
notebooks (EDA, classical models, deep models, analysis) on top of tested modules in
\texttt{code/modules/}; every number in this report is read from the results they save.
